\section{Divisors}
\todo[inline]{Decide if this is viable/necessary for the thesis. We do use this in the 27 lines proof, specifically self-intersection. This could be assumed known. Could stick to Beauville.}

% We will assume that $X$ is a separated Noetherian integral normal scheme throughout this section.

% A \emph{prime divisor} is an integral subscheme $Z \subset X$ of codimension 1.

% \begin{definition}
%     A \emph{Weil Divisor} is a finite formal sum 
%     \[D = m_1 Z_1 + \dots m_r Z_r\]
%     where $m_i \in \mathbb{Z}$ and $Z_i$ is a prime divisor for all $i$. 
% \end{definition}

% The set of divisors will form a group, denoted $\text{Div}(X)$. We say that a divisor $D$ is \emph{effective} if each $m_i \geq 0$, and we write $D \geq 0$. 


% \begin{definition}[\cite{EO25}, Definition 20.4]
%     For a rational function $f \in K(X)^\times$, we define its corresponding divisor by
%     \[\text{div}(f) = \sum_Z \text{ord}_Z(f) Z,\]
%     where the sum runs over all prime divisors. Divisors of the form $\text{div}(f)$ are called \emph{principal divisors}, and they form a subgroup of $\text{Div}(X)$.
% \end{definition}

% \todo[inline]{Consider adding remark about order of $Z$ being well-defined, from discussion in EO or Hartshorne.}

% \begin{definition}[\cite{EO25}, Definition 20.9]
%     The \emph{class group of $X$} is defined as the group of Weil divisors modulo the principal divisors, i.e., 
%     \[\text{Cl}(X) = \text{Div}(X) / \{\text{principal divisors}\}.\]
    
% \end{definition}
